% !PRINTED 500
\weberpage{500}
%<<CONT
if $\varphi\psi_1 = \varphi_1\psi$. If the two functions $\varphi, \psi$ have no common divisor, then $\varphi:\psi$ is called an \gesperrt{irreducible fraction}. Among the various representations of a function $\omega$ there is \gesperrt{one} by an irreducible fraction, and this we shall call the \gesperrt{simplest representation}. In it, numerator and denominator are completely determined by $\omega$ itself, up to a common factor, which may be an arbitrary number in $\Omega$ (Bd. I, \S.\,141).

Such a function $\omega$ we shall call a \gesperrt{functional of the field} $\Omega$. The integral functions and the numbers themselves are contained among these as special cases. For the numbers, \gesperrt{every} representation as a quotient of two numbers in $\Omega$ is to be regarded as a simplest representation.

The four fundamental operations of arithmetic can be applied to the functionals to the same extent as to the numbers, and the totality of all functionals of the field $\Omega$ is therefore likewise a field, which we denote by $\overline{\Omega}$ and shall call the \gesperrt{functional field} of $\Omega$.\footnote{It is natural to designate the functionals of the field $\Omega$, with which one calculates as with the numbers in $\Omega$, outright as \gesperrt{ideal numbers} of the field $\Omega$. This mode of expression would attach itself on the one hand to the ideal factors of \gesperrt{Kummer}, on the other to the ideals introduced by \gesperrt{Dedekind}, which stand to the functionals in a very close relation, to be discussed later. However tempting such a terminology would be in some respects, it proves nevertheless not to be expedient in others, because it would apparently attribute to the functionals --- which are after all only a means to an end, not in themselves an object of our interest --- a greater importance than in reality belongs to them. In introducing the word ``Functional'' we follow a suggestion of \gesperrt{Dedekind}.}

The functional field $\overline{\Omega}$ contains the number field $\Omega$ as a divisor.

A functional whose coefficients are rational numbers is called a \gesperrt{rational functional}. The totality of all rational functionals is the functional field $\overline{R}$ of the field $R$ of the rational numbers, and the field $\overline{R}$ is contained in every algebraic functional field $\overline{\Omega}$.

Every rational functional $A$ can be represented as a quotient of two integral functions in $R$, to which one can also give integral coefficients, by multiplying numerator and denominator by the common denominator of all the coefficients.
